\section*{Глава 1. Анализ предметной области}
\addcontentsline{toc}{section}{Глава 1. Анализ предметной области}
\stepcounter{section}

\subsection{Обзор существующих программных продуктов}

На рынке представлено несколько решений, частично закрывающих потребность
в разговорной практике иностранного языка. Для структурированного сравнения
выделены пять критериев, определяющих пригодность платформы для организации
разговорной языковой практики:

\begin{itemize}
    \item \textbf{Групповой формат встреч} -- разговорная практика эффективнее
          в малой группе, чем в формате один на один: участники слышат разные
          акценты и стили речи, а нагрузка на каждого ниже.
    \item \textbf{Запись участников на встречи} -- необходима для планирования
          мероприятий, контроля числа участников и уведомления организатора.
    \item \textbf{Уровень языка (CEFR) в профиле} -- международная шкала CEFR
          позволяет однозначно указать уровень владения языком и подобрать
          встречу подходящей сложности.
    \item \textbf{Фильтрация встреч по уровню языка} -- без фильтрации участник
          вынужден просматривать все встречи подряд; фильтр по CEFR сокращает
          время поиска и снижает вероятность записи на несоответствующее мероприятие.
    \item \textbf{Языковая специализация} -- платформа, заточенная под языковую
          практику, может предоставлять релевантные инструменты: профиль с уровнем
          языка, тематические клубы, -- чего не даёт универсальный сервис мероприятий.
\end{itemize}

\textbf{HelloTalk} -- мобильное приложение для поиска собеседника по переписке
и голосовым звонкам. Поддерживает уровни языка и интересы в профиле, однако
ориентировано исключительно на общение один на один. Групповой формат и запись
на встречи не предусмотрены.

\textbf{Tandem} -- аналогичное приложение для языкового обмена. Основная механика --
найти носителя языка и переписываться. Клубы и расписание встреч отсутствуют.

\textbf{Meetup} -- платформа для организации мероприятий любой тематики.
Поддерживает создание групп, расписание и запись участников, однако не имеет
языковой специализации: в профиле отсутствуют уровни CEFR, фильтрация по уровню
владения языком не реализована.

Сравнительный анализ аналогов по выделенным критериям представлен в таблице~\ref{tab:analogs}.

\begin{table}[h]
\centering
\caption{Сравнение существующих решений}
\label{tab:analogs}
\begin{tabular}{|p{5.5cm}|c|c|c|c|}
\hline
\textbf{Критерий} & \textbf{HelloTalk} & \textbf{Tandem} & \textbf{Meetup} & \textbf{Наша платформа} \\
\hline
Групповой формат встреч & -- & -- & + & + \\
\hline
Запись участников на встречи & -- & -- & + & + \\
\hline
Уровень языка (CEFR) в профиле & + & + & -- & + \\
\hline
Фильтрация встреч по уровню языка & -- & -- & -- & + \\
\hline
Языковая специализация & + & + & -- & + \\
\hline
\end{tabular}
\end{table}

Анализ показал, что ни одно из рассмотренных решений не сочетает языковую
специализацию и клубный формат с записью на конкретные встречи, что подтверждает
необходимость разработки специализированной платформы.

\subsection{Анализ программных инструментов разработки}

Разработка современного веб-приложения требует выбора инструментов в четырёх
ключевых категориях: fullstack-фреймворк, платформа данных и аутентификации,
CSS-фреймворк и хостинг. В каждой категории существует несколько распространённых
решений; ниже рассмотрены доступные варианты и обоснован итоговый выбор.

\textbf{Fullstack-фреймворк.} Среди актуальных решений выделяются Next.js,
Nuxt.js и SvelteKit. Nuxt.js и SvelteKit ориентированы на экосистемы Vue и Svelte
соответственно, тогда как Next.js построен на React -- наиболее распространённой
библиотеке для построения пользовательских интерфейсов. Next.js 15 предоставляет
встроенный серверный слой (App Router), что позволяет реализовать клиентскую и
серверную части в едином репозитории без отдельного backend-сервера.

\textbf{Платформа данных и аутентификации.} Рассмотрены три варианта: Firebase
(Google), Supabase и самостоятельно развёрнутый PostgreSQL. Firebase использует
документно-ориентированную модель хранения данных, что усложняет реализацию
реляционных связей между сущностями. Самостоятельный PostgreSQL требует отдельной
реализации аутентификации, управления сессиями и хранилища файлов. Supabase
предоставляет реляционную базу данных PostgreSQL, встроенную систему аутентификации
и файловое хранилище в рамках единой облачной платформы, что сокращает объём
инфраструктурного кода.

\textbf{CSS-фреймворк.} Среди распространённых вариантов -- Bootstrap, Tailwind CSS
и CSS Modules. Bootstrap предлагает готовые компоненты с фиксированным визуальным
стилем, что ограничивает гибкость оформления. CSS Modules требуют ручного написания
всех стилей. Tailwind CSS реализует утилитарный подход: стили задаются классами
прямо в разметке, что ускоряет разработку интерфейса без создания отдельных
таблиц стилей.

\textbf{Хостинг.} Рассмотрены Vercel, Netlify и VPS с ручной настройкой (nginx,
SSL, CI/CD). VPS предоставляет максимальный контроль, но требует значительных
затрат на настройку инфраструктуры. Netlify ориентирован преимущественно на
статические сайты. Vercel разработан той же командой, что и Next.js, обеспечивает
автоматический деплой при обновлении репозитория и предоставляется бесплатно
для учебных проектов.

Сравнение вариантов по каждой категории представлено в таблицах
\ref{tab:fw}--\ref{tab:hosting}.

\begin{table}[h]
\centering
\caption{Сравнение fullstack-фреймворков}
\label{tab:fw}
\begin{tabular}{|p{5.5cm}|c|c|c|}
\hline
\textbf{Критерий} & \textbf{Next.js} & \textbf{Nuxt.js} & \textbf{SvelteKit} \\
\hline
Серверный рендеринг (SSR) & + & + & + \\
\hline
Встроенный API без отдельного сервера & + & + & + \\
\hline
Экосистема React & + & -- & -- \\
\hline
Широкое распространение в индустрии & + & -- & -- \\
\hline
\end{tabular}
\end{table}

\begin{table}[h]
\centering
\caption{Сравнение платформ данных и аутентификации}
\label{tab:db}
\begin{tabular}{|p{5.5cm}|c|c|c|}
\hline
\textbf{Критерий} & \textbf{Firebase} & \textbf{Supabase} & \textbf{PostgreSQL} \\
\hline
Реляционная модель данных & -- & + & + \\
\hline
Встроенная аутентификация & + & + & -- \\
\hline
Файловое хранилище & + & + & -- \\
\hline
Открытый исходный код & -- & + & + \\
\hline
\end{tabular}
\end{table}

\begin{table}[h]
\centering
\caption{Сравнение CSS-фреймворков}
\label{tab:css}
\begin{tabular}{|p{5.5cm}|c|c|c|}
\hline
\textbf{Критерий} & \textbf{Bootstrap} & \textbf{Tailwind CSS} & \textbf{CSS Modules} \\
\hline
Гибкость оформления & -- & + & + \\
\hline
Скорость разработки интерфейса & + & + & -- \\
\hline
Отсутствие избыточных стилей & -- & + & + \\
\hline
\end{tabular}
\end{table}

\begin{table}[h]
\centering
\caption{Сравнение платформ хостинга}
\label{tab:hosting}
\begin{tabular}{|p{5.5cm}|c|c|c|}
\hline
\textbf{Критерий} & \textbf{Vercel} & \textbf{Netlify} & \textbf{VPS} \\
\hline
Автодеплой из Git & + & + & -- \\
\hline
Нативная поддержка Next.js & + & -- & -- \\
\hline
Бесплатный тариф & + & + & -- \\
\hline
Не требует настройки сервера & + & + & -- \\
\hline
\end{tabular}
\end{table}

\clearpage

Таким образом, для реализации проекта выбраны Next.js~15, Supabase, Tailwind~CSS
и Vercel -- совокупность инструментов, обеспечивающая fullstack-разработку
в едином репозитории, встроенную аутентификацию и хранилище данных, гибкую
стилизацию и автоматический деплой при минимальных инфраструктурных затратах.

\subsection{Формулировка цели и задач работы}

На основе проведённого анализа существующих продуктов и инструментов разработки
сформулирована цель исследования.

\textbf{Цель работы} -- разработка веб-приложения для поиска языковых клубов
и записи на разговорные встречи по английскому языку.

Для достижения поставленной цели необходимо решить следующие \textbf{задачи}:
\begin{enumerate}
    \item Провести анализ существующих программных продуктов и инструментов разработки;
    \item Спроектировать архитектуру приложения и структуру базы данных;
    \item Реализовать систему аутентификации и разграничения доступа по ролям;
    \item Реализовать CRUD-интерфейс для управления клубами и встречами;
    \item Реализовать запись участников на встречи и фильтрацию по уровню языка и дате;
    \item Обеспечить загрузку аватара и выгрузку списка участников в формате CSV;
    \item Провести тестирование и развернуть приложение на хостинге.
\end{enumerate}
